\documentclass[aspectratio=169,10pt]{beamer}
% Shared typography and drawing styles. Sizes are true TeX points, not scaled.
\usepackage{fontspec}
\setsansfont{LiberationSans-Regular.ttf}[
  Path=fonts/,
  BoldFont=LiberationSans-Bold.ttf,
  ItalicFont=LiberationSans-Italic.ttf,
  BoldItalicFont=LiberationSans-BoldItalic.ttf
]
\usepackage{tikz}
\usetikzlibrary{arrows.meta,calc}

\definecolor{DuneBlue}{HTML}{7CAED5}
\definecolor{DuneOrange}{HTML}{F2682B}
\definecolor{Ink}{HTML}{1C2834}
\definecolor{DeepBlue}{HTML}{304E62}
\definecolor{Rule}{HTML}{CED9E1}
\definecolor{BlueTint}{HTML}{EFF5FA}
\definecolor{OrangeTint}{HTML}{FDF0E9}

\setbeamercolor{normal text}{fg=Ink,bg=white}
\setbeamercolor{structure}{fg=DeepBlue}
\setbeamertemplate{navigation symbols}{}
\setbeamertemplate{headline}{}
\setbeamertemplate{footline}{}
\setbeamersize{text margin left=0.95cm,text margin right=0.95cm}
\setbeamertemplate{background canvas}{%
  \begin{tikzpicture}[remember picture,overlay]
    \fill[white] (current page.south west) rectangle (current page.north east);
    \fill[DuneBlue] (current page.north west) rectangle
      ($(current page.north east)+(0,-0.05)$);
    \node[anchor=north west,inner sep=0pt,text=DeepBlue,
      font=\fontsize{6.5}{8}\selectfont]
      at ($(current page.north west)+(0.95,-0.38)$)
      {DUNE FAR DETECTOR\quad /\quad SURF};
    \node[anchor=south east,inner sep=0pt,text=DeepBlue,
      font=\fontsize{6.5}{8}\selectfont]
      at ($(current page.south east)+(-0.95,0.34)$)
      {\insertframenumber\,/\,\inserttotalframenumber};
  \end{tikzpicture}%
}
\hypersetup{
  pdftitle={SURF--DUNE facility interface},
  pdfauthor={DUNE Far Detector},
  colorlinks=true,urlcolor=DeepBlue
}
\hyphenpenalty=10000
\exhyphenpenalty=10000
\tracinglostchars=3
\raggedright

\tikzset{
  heading/.style={anchor=north west,inner sep=0pt,text=Ink,align=left,
    font=\fontsize{18}{21}\selectfont\bfseries},
  subheading/.style={anchor=north west,inner sep=0pt,text=Ink,align=left,
    font=\fontsize{11.5}{14}\selectfont\bfseries},
  body/.style={anchor=north west,inner sep=0pt,text=Ink,align=left,
    font=\fontsize{10.5}{13.5}\selectfont},
  small/.style={anchor=north west,inner sep=0pt,text=DeepBlue,align=left,
    font=\fontsize{9}{11.5}\selectfont},
  label/.style={anchor=north west,inner sep=0pt,text=DeepBlue,align=left,
    font=\fontsize{7.5}{9.5}\selectfont\bfseries},
  zone/.style={draw=Rule,fill=white,rounded corners=2pt,line width=0.55pt},
  point/.style={draw=DeepBlue,line width=0.65pt,fill=white,
    rounded corners=1.5pt,inner sep=3pt,align=center,
    font=\fontsize{9.2}{11.2}\selectfont,text=Ink},
  flow/.style={-{Latex[length=1.8mm,width=1.3mm]},draw=DeepBlue,line width=0.75pt},
  protectionflow/.style={flow,draw=Ink,line width=1pt},
  statusflow/.style={flow,dashed,line width=0.65pt},
  edgelabel/.style={fill=white,inner xsep=3pt,inner ysep=1.5pt,
    text=DeepBlue,font=\fontsize{7.4}{9}\selectfont,align=center}
}
\newenvironment{slidecanvas}{%
  \begin{tikzpicture}[remember picture,overlay,
    shift={($(current page.north west)+(0.95,-1.05)$)},x=1cm,y=1cm]
}{\end{tikzpicture}}
\newcommand{\SlideTitle}[2]{%
  \node[heading] at (0,0) {#1};
  \node[small,text width=14.1cm] at (0,-0.80) {#2};
}
\newcommand{\TableRow}[3]{%
  \node[body,text width=3.8cm] at (0,#1) {#2};
  \node[body,text width=10.05cm] at (4.05,#1) {#3};
}
\newcommand{\NodeDetail}[1]{\\[1.3pt]{\fontsize{7.8}{9.2}\selectfont #1}}


\begin{document}

\begin{frame}[plain]
  \begin{slidecanvas}
    \SlideTitle{How facility signals reach DUNE}
      {Proposed paths for monitoring and detector protection.}
    % Adapted from the earlier three-zone reference architecture (diagram 01).
% Positions are centimetres from the common slide-content origin.
% The zones end above the three notes; all connectors use explicit node anchors.
\draw[zone] (0,-1.65) rectangle (3.05,-5.90);
\draw[zone] (3.32,-1.65) rectangle (6.15,-5.90);
\draw[zone] (6.42,-1.65) rectangle (14.10,-5.90);
\fill[BlueTint] (0.01,-1.66) rectangle (3.04,-2.15);
\fill[OrangeTint] (3.33,-1.66) rectangle (6.14,-2.15);
\fill[BlueTint] (6.43,-1.66) rectangle (14.09,-2.15);
\node[label,font=\fontsize{8.3}{10}\selectfont\bfseries]
  at (0.17,-1.80) {SURF / sources};
\node[label,font=\fontsize{7.2}{9}\selectfont\bfseries]
  at (3.49,-1.80) {Controlled boundary};
\node[label,font=\fontsize{8.3}{10}\selectfont\bfseries]
  at (6.59,-1.80) {DUNE experiment};

\node[label] at (0.17,-2.39) {MONITORING};
\node[point,minimum width=2.62cm,text width=2.32cm,minimum height=1.05cm]
  (bms) at (1.525,-3.15) {BMS / Metasys};
\node[point,minimum width=2.38cm,text width=2.08cm,minimum height=1.05cm]
  (boundary) at (4.735,-3.15) {Read-only\\exchange};
\node[point,minimum width=2.04cm,text width=1.74cm,minimum height=1.05cm]
  (reader) at (7.685,-3.15) {Facility\\reader};
\node[point,minimum width=2.04cm,text width=1.74cm,minimum height=1.05cm]
  (slow) at (10.195,-3.15) {Slow\\Controls};
\node[point,minimum width=2.27cm,text width=1.97cm,minimum height=1.05cm]
  (operators) at (12.705,-3.15) {Operators};
\draw[flow] (bms.east) -- (boundary.west);
\draw[flow] (boundary.east) -- (reader.west);
\draw[flow] (reader.east) -- (slow.west);
\draw[flow] (slow.east) -- (operators.west);

\node[label,text=Ink] at (0.17,-4.20) {PROTECTION};
\node[point,minimum width=2.62cm,text width=2.32cm,minimum height=1.15cm]
  (source) at (1.525,-5.10) {Approved\\source state};
\node[point,dashed,minimum width=2.38cm,text width=2.08cm,minimum height=1.15cm]
  (dpsinput) at (4.735,-5.10) {Supervised\\DPS input};
\node[point,draw=Ink,fill=Ink,text=white,
  minimum width=2.04cm,text width=1.74cm,minimum height=1.15cm]
  (dps) at (7.685,-5.10) {DPS\NodeDetail{Action / latch}};
\draw[DuneOrange,line width=1.4pt] (dps.south west) -- (dps.south east);
\node[point,minimum width=4.63cm,text width=4.33cm,minimum height=1.15cm]
  (functions) at (11.45,-5.10)
  {Protected detector functions\NodeDetail{Local protection stays separate}};
\draw[protectionflow] (source.east) -- (dpsinput.west);
\draw[protectionflow] (dpsinput.east) -- (dps.west);
\draw[protectionflow] (dps.east) -- (functions.west);
\draw[statusflow] (dps.north) -- (7.685,-4.00) --
  node[edgelabel] {DPS status} (10.195,-4.00) -- (slow.south);

\foreach \x/\accent in {0/DeepBlue,4.95/DuneOrange,9.90/Ink}
  \draw[\accent,line width=1.2pt] (\x,-6.25) -- (\x,-6.92);
\node[label,font=\fontsize{8.6}{10.5}\selectfont\bfseries]
  at (0.16,-6.25) {Functional view};
\node[small,font=\fontsize{8}{10}\selectfont,text width=4.05cm]
  at (0.16,-6.62) {Protection link to be defined.};
\node[label,font=\fontsize{8.6}{10.5}\selectfont\bfseries]
  at (5.11,-6.25) {Lost monitoring};
\node[small,font=\fontsize{8}{10}\selectfont,text width=4.05cm]
  at (5.11,-6.62) {Missing or old data is visible.};
\node[label,font=\fontsize{8.6}{10.5}\selectfont\bfseries]
  at (10.06,-6.25) {Recovery};
\node[small,font=\fontsize{8}{10}\selectfont,text width=4.04cm]
  at (10.06,-6.62) {Restart stays with its owner.};

  \end{slidecanvas}
\end{frame}

\begin{frame}[plain]
  \begin{slidecanvas}
    \SlideTitle{Alarm flow and ownership}
      {Information moves through the experiment; authority stays with the owner.}
    % Adapted from diagram 02: parallel source / DPS authority, then supervision.
% Solid arrows carry condition or operational information; dashed arrows are
% status copies, not control or reset paths. No numbered authority levels.
\node[point,fill=BlueTint,text width=4.3cm,minimum width=4.6cm,minimum height=0.95cm]
  (field) at (7.05,-1.98) {Source systems\NodeDetail{Local sensing and protection}};
\node[point,text width=4.9cm,minimum width=5.2cm,minimum height=0.95cm]
  (facility) at (3.00,-3.30) {Facility or source owner\NodeDetail{Original alarm, response and reset}};
\node[point,fill=Ink,draw=Ink,text=white,text width=4.9cm,
  minimum width=5.2cm,minimum height=0.95cm]
  (protection) at (11.10,-3.30) {Detector protection (DPS)\NodeDetail{Approved action, latch and reset}};
\draw[DuneOrange,line width=1.4pt] (protection.south west) -- (protection.south east);
\node[point,fill=BlueTint,text width=4.9cm,minimum width=5.2cm,minimum height=0.95cm]
  (supervision) at (7.05,-4.66) {Slow Controls\NodeDetail{Time, quality and detector context}};
\node[point,text width=4.9cm,minimum width=5.2cm,minimum height=0.95cm]
  (runcontrol) at (3.00,-6.15) {Run Control\NodeDetail{Run state and orderly stop}};
\node[point,text width=4.9cm,minimum width=5.2cm,minimum height=0.95cm]
  (operations) at (11.10,-6.15) {Operations\NodeDetail{Notification to the responsible team}};

\draw[flow] ($(field.south)+(-1.10,0)$) -- (5.95,-2.62) --
  node[edgelabel] {Facility condition} (3.00,-2.62) -- (facility.north);
\draw[protectionflow] ($(field.south)+(1.10,0)$) -- (8.15,-2.62) --
  node[edgelabel,text=Ink] {Approved state} (11.10,-2.62) -- (protection.north);
\draw[statusflow] (facility.south) -- (3.00,-4.00) --
  node[edgelabel] {Alarm copy} (6.10,-4.00) -- ($(supervision.north)+(-0.95,0)$);
\draw[statusflow] (protection.south) -- (11.10,-4.00) --
  node[edgelabel] {DPS status} (8.00,-4.00) -- ($(supervision.north)+(0.95,0)$);
\draw[flow] ($(supervision.south)+(-0.95,0)$) -- (6.10,-5.44) --
  node[edgelabel] {Availability / stop} (3.00,-5.44) -- (runcontrol.north);
\draw[flow] ($(supervision.south)+(0.95,0)$) -- (8.00,-5.44) --
  node[edgelabel] {Notification} (11.10,-5.44) -- (operations.north);
\draw[flow] (runcontrol.east) -- node[edgelabel] {Run impact} (operations.west);

\node[small,font=\fontsize{8.6}{10.8}\selectfont,text width=14.1cm]
  at (0,-7.00)
  {DUNE acknowledgment applies to the displayed copy, not the source alarm or DPS reset.};

  \end{slidecanvas}
\end{frame}

\begin{frame}[plain]
  \begin{slidecanvas}
    \SlideTitle{Facility conditions}
      {The source definition and agreed cause and effect determine the response.}
    \node[label] at (0,-1.70) {CONDITION};
    \node[label] at (4.05,-1.70) {POSSIBLE DUNE RESPONSE};
    \foreach \position in {-2.05,-2.76,-3.47,-4.18,-4.89,-5.60,-6.31}
      \draw[Rule,line width=0.45pt] (0,\position) -- (14.1,\position);
    \TableRow{-2.27}{Power / reserve}{An agreed request for an orderly DAQ stop.}
    \TableRow{-2.98}{Network / BMS loss}{Stale or unknown data; the loss is visible.}
    \TableRow{-3.69}{Fire / low oxygen}{Source response; detector action only where approved.}
    \TableRow{-4.40}{Cryogenic / LAr leak}{Response follows an explicitly defined source condition.}
    \TableRow{-5.11}{Cooling / water}{Local protection; agreed DPS action where applicable.}
    \TableRow{-5.82}{Weather / access}{SURF operating restrictions guide the response.}
    \node[small,text width=14.1cm] at (0,-6.75)
      {Loss of monitoring is a design case for the independent protection path.};
  \end{slidecanvas}
\end{frame}

\begin{frame}[plain]
  \begin{slidecanvas}
    \SlideTitle{Proposed data exchange}
      {Protocol support and source ownership remain to be confirmed.}
    \node[label] at (0,-1.70) {INTERFACE};
    \node[label] at (3.15,-1.70) {CANDIDATE};
    \node[label] at (6.0,-1.70) {INFORMATION FLOW};
    \foreach \position in {-2.05,-3.13,-4.21,-5.29}
      \draw[Rule,line width=0.45pt] (0,\position) -- (14.1,\position);
    \node[body] at (0,-2.43) {BMS readout};
    \node[body] at (3.15,-2.43) {BACnet/IP};
    \node[body,text width=8.1cm] at (6.0,-2.43) {Reader queries selected source points.};
    \node[body] at (0,-3.51) {DUNE view};
    \node[body] at (3.15,-3.51) {OPC UA};
    \node[body,text width=8.1cm] at (6.0,-3.51) {Reader publishes to Slow Controls.};
    \node[body] at (0,-4.59) {Equipment health};
    \node[body] at (3.15,-4.59) {SNMPv3};
    \node[body,text width=8.1cm] at (6.0,-4.59) {Supported power and network telemetry.};
    \node[small,text width=14.1cm] at (0,-5.84)
      {Read access exchanges requests and replies without granting control authority.};
    \node[small,text width=14.1cm] at (0,-6.34)
      {BACnet/SC is an option where supported. The DPS interface is defined separately.};
  \end{slidecanvas}
\end{frame}

\begin{frame}[plain]
  \begin{slidecanvas}
    \SlideTitle{From concept to agreed interface}
      {Three areas for discussion with the source and experiment teams.}
    \foreach \position/\number in {-1.95/1,-3.50/2,-5.05/3}{
      \node[anchor=north west,inner sep=0pt,text=DeepBlue,
        font=\fontsize{16}{19}\selectfont] at (0,\position) {\number};
      \draw[Rule,line width=0.45pt] (1.0,\position-1.13) -- (14.1,\position-1.13);
    }
    \node[subheading] at (1.0,-1.95) {Point meaning and ownership};
    \node[body,text width=13.1cm] at (1.0,-2.48)
      {Each point has a defined source, state, quality and owner.};
    \node[subheading] at (1.0,-3.50) {Data path and timing};
    \node[body,text width=13.1cm] at (1.0,-4.03)
      {The interface describes read access, update rate and data age.};
    \node[subheading] at (1.0,-5.05) {Response and recovery};
    \node[body,text width=13.1cm] at (1.0,-5.58)
      {Checks cover alarm routing, loss, approved protection action and reset.};
  \end{slidecanvas}
\end{frame}

\end{document}
